\documentclass[a4paper,12pt,twoside]{memoir}

% Castellano
\usepackage[spanish,es-tabla]{babel}
\selectlanguage{spanish}
\usepackage[utf8]{inputenc}
\usepackage[T1]{fontenc}
\usepackage{lmodern} % Scalable font
\usepackage{microtype}
\usepackage{placeins}

\RequirePackage{booktabs}
\RequirePackage[table]{xcolor}
\RequirePackage{xtab}
\RequirePackage{multirow}

% Links
\PassOptionsToPackage{hyphens}{url}\usepackage[colorlinks]{hyperref}
\hypersetup{
	allcolors = {red}
}

% Ecuaciones
\usepackage{amsmath}

% Rutas de fichero / paquete
\newcommand{\ruta}[1]{{\sffamily #1}}

% Párrafos
\nonzeroparskip

% Huérfanas y viudas
\widowpenalty100000
\clubpenalty100000

% Imágenes

% Comando para insertar una imagen en un lugar concreto.
% Los parámetros son:
% 1 --> Ruta absoluta/relativa de la figura
% 2 --> Texto a pie de figura
% 3 --> Tamaño en tanto por uno relativo al ancho de página
\usepackage{graphicx}
\newcommand{\imagen}[3]{
	\begin{figure}[!h]
		\centering
		\includegraphics[width=#3\textwidth]{#1}
		\caption{#2}\label{fig:#1}
	\end{figure}
	\FloatBarrier
}

% Comando para insertar una imagen sin posición.
% Los parámetros son:
% 1 --> Ruta absoluta/relativa de la figura
% 2 --> Texto a pie de figura
% 3 --> Tamaño en tanto por uno relativo al ancho de página
\newcommand{\imagenflotante}[3]{
	\begin{figure}
		\centering
		\includegraphics[width=#3\textwidth]{#1}
		\caption{#2}\label{fig:#1}
	\end{figure}
}

% El comando \figura nos permite insertar figuras comodamente, y utilizando
% siempre el mismo formato. Los parametros son:
% 1 --> Porcentaje del ancho de página que ocupará la figura (de 0 a 1)
% 2 --> Fichero de la imagen
% 3 --> Texto a pie de imagen
% 4 --> Etiqueta (label) para referencias
% 5 --> Opciones que queramos pasarle al \includegraphics
% 6 --> Opciones de posicionamiento a pasarle a \begin{figure}
\newcommand{\figuraConPosicion}[6]{%
  \setlength{\anchoFloat}{#1\textwidth}%
  \addtolength{\anchoFloat}{-4\fboxsep}%
  \setlength{\anchoFigura}{\anchoFloat}%
  \begin{figure}[#6]
    \begin{center}%
      \Ovalbox{%
        \begin{minipage}{\anchoFloat}%
          \begin{center}%
            \includegraphics[width=\anchoFigura,#5]{#2}%
            \caption{#3}%
            \label{#4}%
          \end{center}%
        \end{minipage}
      }%
    \end{center}%
  \end{figure}%
}

%
% Comando para incluir imágenes en formato apaisado (sin marco).
\newcommand{\figuraApaisadaSinMarco}[5]{%
  \begin{figure}%
    \begin{center}%
    \includegraphics[angle=90,height=#1\textheight,#5]{#2}%
    \caption{#3}%
    \label{#4}%
    \end{center}%
  \end{figure}%
}
% Para las tablas
\newcommand{\otoprule}{\midrule [\heavyrulewidth]}
%
% Nuevo comando para tablas pequeñas (menos de una página).
\newcommand{\tablaSmall}[5]{%
 \begin{table}
  \begin{center}
   \rowcolors {2}{gray!35}{}
   \begin{tabular}{#2}
    \toprule
    #4
    \otoprule
    #5
    \bottomrule
   \end{tabular}
   \caption{#1}
   \label{tabla:#3}
  \end{center}
 \end{table}
}

%
% Nuevo comando para tablas pequeñas (menos de una página).
\newcommand{\tablaSmallSinColores}[5]{%
 \begin{table}[H]
  \begin{center}
   \begin{tabular}{#2}
    \toprule
    #4
    \otoprule
    #5
    \bottomrule
   \end{tabular}
   \caption{#1}
   \label{tabla:#3}
  \end{center}
 \end{table}
}

\newcommand{\tablaApaisadaSmall}[5]{%
\begin{landscape}
  \begin{table}
   \begin{center}
    \rowcolors {2}{gray!35}{}
    \begin{tabular}{#2}
     \toprule
     #4
     \otoprule
     #5
     \bottomrule
    \end{tabular}
    \caption{#1}
    \label{tabla:#3}
   \end{center}
  \end{table}
\end{landscape}
}

%
% Nuevo comando para tablas grandes con cabecera y filas alternas coloreadas en gris.
\newcommand{\tabla}[6]{%
  \begin{center}
    \tablefirsthead{
      \toprule
      #5
      \otoprule
    }
    \tablehead{
      \multicolumn{#3}{l}{\small\sl continúa desde la página anterior}\\
      \toprule
      #5
      \otoprule
    }
    \tabletail{
      \hline
      \multicolumn{#3}{r}{\small\sl continúa en la página siguiente}\\
    }
    \tablelasttail{
      \hline
    }
    \bottomcaption{#1}
    \rowcolors {2}{gray!35}{}
    \begin{xtabular}{#2}
      #6
      \bottomrule
    \end{xtabular}
    \label{tabla:#4}
  \end{center}
}

%
% Nuevo comando para tablas grandes con cabecera.
\newcommand{\tablaSinColores}[6]{%
  \begin{center}
    \tablefirsthead{
      \toprule
      #5
      \otoprule
    }
    \tablehead{
      \multicolumn{#3}{l}{\small\sl continúa desde la página anterior}\\
      \toprule
      #5
      \otoprule
    }
    \tabletail{
      \hline
      \multicolumn{#3}{r}{\small\sl continúa en la página siguiente}\\
    }
    \tablelasttail{
      \hline
    }
    \bottomcaption{#1}
    \begin{xtabular}{#2}
      #6
      \bottomrule
    \end{xtabular}
    \label{tabla:#4}
  \end{center}
}

%
% Nuevo comando para tablas grandes sin cabecera.
\newcommand{\tablaSinCabecera}[5]{%
  \begin{center}
    \tablefirsthead{
      \toprule
    }
    \tablehead{
      \multicolumn{#3}{l}{\small\sl continúa desde la página anterior}\\
      \hline
    }
    \tabletail{
      \hline
      \multicolumn{#3}{r}{\small\sl continúa en la página siguiente}\\
    }
    \tablelasttail{
      \hline
    }
    \bottomcaption{#1}
  \begin{xtabular}{#2}
    #5
   \bottomrule
  \end{xtabular}
  \label{tabla:#4}
  \end{center}
}



\definecolor{cgoLight}{HTML}{EEEEEE}
\definecolor{cgoExtralight}{HTML}{FFFFFF}

%
% Nuevo comando para tablas grandes sin cabecera.
\newcommand{\tablaSinCabeceraConBandas}[5]{%
  \begin{center}
    \tablefirsthead{
      \toprule
    }
    \tablehead{
      \multicolumn{#3}{l}{\small\sl continúa desde la página anterior}\\
      \hline
    }
    \tabletail{
      \hline
      \multicolumn{#3}{r}{\small\sl continúa en la página siguiente}\\
    }
    \tablelasttail{
      \hline
    }
    \bottomcaption{#1}
    \rowcolors[]{1}{cgoExtralight}{cgoLight}

  \begin{xtabular}{#2}
    #5
   \bottomrule
  \end{xtabular}
  \label{tabla:#4}
  \end{center}
}



\graphicspath{ {./img/} }

% Capítulos
\chapterstyle{bianchi}
\newcommand{\capitulo}[2]{
	\setcounter{chapter}{#1}
	\setcounter{section}{0}
	\setcounter{figure}{0}
	\setcounter{table}{0}
	\chapter*{\thechapter.\enskip #2}
	\addcontentsline{toc}{chapter}{\thechapter.\enskip #2}
	\markboth{#2}{#2}
}

% Apéndices
\renewcommand{\appendixname}{Apéndice}
\renewcommand*\cftappendixname{\appendixname}

\newcommand{\apendice}[1]{
	%\renewcommand{\thechapter}{A}
	\chapter{#1}
}

\renewcommand*\cftappendixname{\appendixname\ }

% Formato de portada
\makeatletter
\usepackage{xcolor}
\newcommand{\tutor}[1]{\def\@tutor{#1}}
\newcommand{\course}[1]{\def\@course{#1}}
\definecolor{cpardoBox}{HTML}{E6E6FF}
\def\maketitle{
  \null
  \thispagestyle{empty}
  % Cabecera ----------------
\noindent\includegraphics[width=\textwidth]{cabecera}\vspace{1cm}%
  \vfill
  % Título proyecto y escudo informática ----------------
  \colorbox{cpardoBox}{%
    \begin{minipage}{.8\textwidth}
      \vspace{.5cm}\Large
      \begin{center}
      \textbf{TFG del Grado en Ingeniería Informática}\vspace{.6cm}\\
      \textbf{\LARGE\@title{}}
      \end{center}
      \vspace{.2cm}
    \end{minipage}

  }%
  \hfill\begin{minipage}{.20\textwidth}
    \includegraphics[width=\textwidth]{escudoInfor}
  \end{minipage}
  \vfill
  % Datos de alumno, curso y tutores ------------------
  \begin{center}%
  {%
    \noindent\LARGE
    Presentado por \@author{}\\ 
    en Universidad de Burgos --- \@date{}\\
    Tutor: \@tutor{}\\
  }%
  \end{center}%
  \null
  \cleardoublepage
  }
\makeatother

\newcommand{\nombre}{Alejandro de Castro Arnaiz} %%% cambio de comando

% Datos de portada
\title{Open Data in Smart Cities}
\author{\nombre}
\tutor{Bruno Baruque Zanon y Héctor Cogollos Adrián}
\date{\today}

\begin{document}

\maketitle


\newpage\null\thispagestyle{empty}%\newpage


%%%%%%%%%%%%%%%%%%%%%%%%%%%%%%%%%%%%%%%%%%%%%%%%%%%%%%%%%%%%%%%%%%%%%%%%%%%%%%%%%%%%%%%%
\thispagestyle{empty}


\noindent\includegraphics[width=\textwidth]{cabecera}\vspace{1cm}

\noindent D. nombre tutor, profesor del departamento de nombre departamento, área de nombre área.

\noindent Expone:

\noindent Que el alumno D. \nombre, con DNI 71792120N, ha realizado el Trabajo final de Grado en Ingeniería Informática titulado Open Data in Smart Cities de TFG. 

\noindent Y que dicho trabajo ha sido realizado por el alumno bajo la dirección del que suscribe, en virtud de lo cual se autoriza su presentación y defensa.

\begin{center} %\large
En Burgos, {\large \today}
\end{center}

\vfill\vfill\vfill

% Author and supervisor
\begin{minipage}{0.45\textwidth}
\begin{flushleft} %\large
Vº. Bº. del Tutor:\\[2cm]
D. Bruno Baruque Zanon
\end{flushleft}
\end{minipage}
\hfill
\begin{minipage}{0.45\textwidth}
\begin{flushleft} %\large
Vº. Bº. del co-tutor:\\[2cm]
D. Héctor Cogollos Adrián
\end{flushleft}
\end{minipage}
\hfill

\vfill

% para casos con solo un tutor comentar lo anterior
% y descomentar lo siguiente
%Vº. Bº. del Tutor:\\[2cm]
%D. nombre tutor


\newpage\null\thispagestyle{empty}%\newpage




\frontmatter

% Abstract en castellano
\renewcommand*\abstractname{Resumen}
\begin{abstract}
\begin{comment}
En este primer apartado se hace una \textbf{breve} presentación del tema que se aborda en el proyecto.
\end{comment}
Durante los últimos años, el acceso a datos en tiempo real gracias a sensores ha incrementado exponencialmente, dejando así, una versión nueva del Internet de las cosas (\textit{IOT}).
Esto permite que estén disponibles datos como aforos, horarios, precios y disponibilidad entre otros muchos datos de diferentes edificios y servicios en las ciudades. La gran cantidad de organismos que ofrecen datos abiertos provoca que haya multitud de formatos y estructuras de datos diferentes, lo que dificulta unificarlas en aplicaciones funcionales. Es por ello que hace unos años se lanzó \textit{FIWARE}, una plataforma europea que cuenta con el apoyo con de las principales empresas tecnológicas y que propone un estándar para dichos datos.
El proyecto consiste en el desarrollo de una aplicación que permita al usuario una experiencia más cómoda en una ciudad, mediante el tratamiento de datos abiertos incorporando \textit{FIWARE}.
Se ofrece al usuario diversas funcionalidades con los datos obtenidos de su ciudad (aforos, localizaciones, teléfonos, horarios...), como la consulta, registro de datos periódicamente, entre otras funcionalidades. Se ha desarrollado de manera interactiva, dejando así al usuario la capacidad de añadir nuevos conjuntos de datos.
\end{abstract}

\renewcommand*\abstractname{Descriptores}
\begin{abstract}
\begin{comment}
Palabras separadas por comas que identifiquen el contenido del proyecto Ej: servidor web, buscador de vuelos, android \ldots
\end{comment}
fiware, iot, open data, ngsi, datos en tiempo real, conjunto de datos, datos abiertos
\end{abstract}

\clearpage

% Abstract en inglés
\renewcommand*\abstractname{Abstract}
\begin{abstract}
A \textbf{brief} presentation of the topic addressed in the project. COMPLETAR CUANDO ESTÉ CORREGIDO EN ESPAÑOL
\end{abstract}

\renewcommand*\abstractname{Keywords}
\begin{abstract}
keywords separated by commas.
\end{abstract}

\clearpage

% Indices
\tableofcontents

\clearpage

\listoffigures

\clearpage

\listoftables
\clearpage

\mainmatter
\capitulo{1}{Introducción}
\begin{comment}
Descripción del contenido del trabajo y del estructura de la memoria y del resto de materiales entregados.
\end{comment}

En la última década, la demanda de acceso instantáneo a datos en tiempo real ha crecido exponencialmente. Esta necesidad se debe en parte a la proliferación de dispositivos conectados y al avance de la tecnología de la información y la comunicación (TIC). Los seres humanos, tanto a nivel individual como colectivo, requieren información rápida y precisa para tomar decisiones informadas, mejorar la eficiencia operativa y optimizar la gestión de recursos \cite{khajeh2020, patel2016, gubbi2013}.

Para acceder a la información en tiempo real, es fundamental el uso de dispositivos equipados con sensores y conexión a Internet. Estos dispositivos, que pueden variar desde teléfonos inteligentes y relojes inteligentes hasta sensores ambientales y cámaras de seguridad, son capaces de recopilar, procesar y transmitir datos de manera continua. Este flujo constante de datos permite a los sistemas y aplicaciones reaccionar instantáneamente a las condiciones cambiantes y proporcionar información actualizada a los usuarios.

El concepto de Internet de las Cosas (\textit{IoT}, por sus siglas en inglés) se refiere a la interconexión de dispositivos y objetos cotidianos a través de Internet, permitiéndoles enviar y recibir datos. Según Gubbi et al. (2013), IoT es ''una visión en la que 'cosas' (objetos físicos) están integradas con electrónica, software, sensores y conectividad de red para recopilar e intercambiar datos'' \cite{gubbi2013}. Esto crea un ecosistema en el que los datos pueden ser utilizados para mejorar la eficiencia, productividad y calidad de vida.

En el contexto de las \textit{Smart Cities}, el \textit{IoT} juega un papel crucial. Los sensores y dispositivos IoT distribuidos por toda la ciudad pueden monitorizar el estado de diferentes infraestructuras, precios, disponibilidad, y otros parámetros esenciales en tiempo real.

Este trabajo propone una aplicación que facilita al usuario el acceso a todos estos datos, permitiendo así, a cualquier persona acceder a los datos de la ciudad que desee en tiempo real, de manera rápida y precisa. 

Existen otras aplicaciones que permiten el acceso a dichos datos, pero muchas de ellas suelen estar enfocadas en un conjunto de datos concreto, como pueden ser los \textit{parkings}, en una ciudad específica, o simplemente son \textit{APIs} de datos que ofrecen los datos en crudo mediante ficheros \textit{JSON}, \textit{CSV} u otros formatos determinados.
En esta aplicación, el usuario puede acceder al conjunto de datos que desee, de la ciudad que desee, y tener una visión gráfica y personalizada de los distintos conjuntos de datos. Además, esta aplicación tiene la ventaja de estar desarrollada de manera web, esto permite el acceso rápido a cualquier usuario y en cualquier dispositivo que tenga acceso a un navegador web, sin necesidad de descargarse ninguna aplicación.

\section{Estructura de la memoria}\label{estructura-de-la-memoria}

La memoria sigue la siguiente estructura:

\begin{itemize}
\tightlist
\item
  \textbf{Introducción:} breve introducción al problema a resolver y la
  solución propuesta en este trabajo. Estructura de la memoria y listado de materiales
  entregados.
\item
  \textbf{Objetivos del proyecto:} objetivos que
  persigue el desarrollo del proyecto. Tanto los objetivos marcados por los requisitos del software y los objetivos de carácter técnico.
\item
  \textbf{Conceptos teóricos:} explicación de los conceptos teóricos más importantes para la comprensión del proyecto.
\item
  \textbf{Técnicas y herramientas:} técnicas y herramientas utilizadas durante las diferentes etapas del proyecto y una explicación y/o comparación con otras que finalmente no se han utilizado.
\item
  \textbf{Aspectos relevantes del desarrollo:} aspectos más importantes del proyecto, como sus etapas, metodología, el proceso de aprendizaje...
\item
  \textbf{Trabajos relacionados:} 
\item
  \textbf{Conclusiones y líneas de trabajo futuras:} conclusiones del proyecto, conclusiones técnicas, conclusiones personales y futuras mejoras.
\end{itemize}

La memoria contiene los siguientes anexos:

\begin{itemize}
\tightlist
\item
  \textbf{Plan del proyecto software:} planificación temporal del proyecto, y viabilidad económica y legal del mismo.
\item
  \textbf{Especificación de requisitos del software:} se detallan los objetivos, los requisitos funcionales, así como los casos de uso de la aplicación.
\item
  \textbf{Especificación de diseño:} se detallan las guías del diseño empleadas para el desarrollo del proyecto, tales como diagrama de secuencias, modelado de datos o diagrama de interfaces.
\item
  \textbf{Manual del programador:} EN PROCESO
\item
  \textbf{Manual de usuario:} guía para el usuario donde se incluyen los requisitos para el uso de la aplicación y un manual para su uso.
\item
  \textbf{ODS:} AÚN POR HACER
\end{itemize}






\capitulo{2}{Objetivos del proyecto}
\begin{comment}
Este apartado explica de forma precisa y concisa cuales son los objetivos que se persiguen con la realización del proyecto. Se puede distinguir entre los objetivos marcados por los requisitos del software a construir y los objetivos de carácter técnico que plantea a la hora de llevar a la práctica el proyecto.
\end{comment}

El objetivo principal del proyecto es el desarrollo de una aplicación de datos abiertos en \textit{Smart Cities} utilizando la plataforma FIWARE \cite{fiware}. 

\section{Objetivos marcados por los requisitos del \textit{software}}\label{objetivos-generales}

\begin{itemize}
\tightlist
\item
  Permitir a cualquier usuario el acceso a un conjunto de datos de una ciudad de manera rápida y precisa.
\item
  Ofrecer al usuario una experiencia sencilla y cómoda, enfocada a cualquier tipo de público. También se busca que esta experiencia sea personalizada, pudiendo personalizar cada conjunto por cada usuario.
\item
  Facilitar al usuario la incorporación de nuevos conjuntos de datos, disponibles para todos los usuarios.
\item
  Permitir al usuario guardar un registro de los conjuntos de datos que desee para un posterior análisis.
\end{itemize}



\section{Objetivos de carácter técnico}\label{objetivos-personales}

\begin{itemize}
\tightlist

\item Investigar sobre la implantación de \textit{FIWARE} y ponerlo en marcha.
\item Aplicar y ampliar los conocimientos adquiridos durante el grado.
\item Aprender sobre el desarrollo de aplicaciones web.
\item Ampliar y consolidar conocimientos sobre análisis y tratamiento de datos.
\item Usar GitHub para el control de versiones.
\item Implementar una base de datos en la aplicación para guardar los diferentes datos necesarios.
\item Ampliar los conocimientos de la herramienta de desarrollo de documentos \LaTeX{}.
\item Usar una metodología ágil \textit{Scrum} para el desarrollo del proyecto.
\item Aplicar y ampliar los conocimientos sobre \textit{Python}, desarrollo web, bases de datos, gestión de proyectos, estructuras de datos y análisis de \textit{software} adquiridos en el grado.


\end{itemize}
\capitulo{3}{Conceptos teóricos}

En aquellos proyectos que necesiten para su comprensión y desarrollo de unos conceptos teóricos de una determinada materia o de un determinado dominio de conocimiento, debe existir un apartado que sintetice dichos conceptos.

Algunos conceptos teóricos de \LaTeX{} \footnote{Créditos a los proyectos de Álvaro López Cantero: Configurador de Presupuestos y Roberto Izquierdo Amo: PLQuiz}.

FIWARE es una plataforma de código abierto que facilita el desarrollo de soluciones inteligentes, gemelos digitales y espacios de datos en varios dominios de la transformación digital. FIWARE se basa en el uso de APIs estándar para gestionar la información de contexto, lo que permite actualizar y acceder a los datos de forma eficiente y segura. FIWARE cuenta con el apoyo de la Unión Europea y de una comunidad global de desarrolladores, proveedores de servicios, empresas y organizaciones, las cuales en 2012 anunciaron el proyecto de desarrollar estándares de FIWARE para smart cities.
Esta plataforma permite dirigir las aplicaciones Open Data, Data Lakes y gestores de contexto.
Fiware, se apoya en software de código abierto, del cual ya disfrutan 350 ciudades.



\section{Secciones}

Las secciones se incluyen con el comando section.

\subsection{Subsecciones}

Además de secciones tenemos subsecciones.

\subsubsection{Subsubsecciones}

Y subsecciones. 


\section{Referencias}

Las referencias se incluyen en el texto usando cite~\cite{wiki:latex}. Para citar webs, artículos o libros~\cite{koza92}, si se desean citar más de uno en el mismo lugar~\cite{bortolot2005, koza92}.


\section{Imágenes}

Se pueden incluir imágenes con los comandos standard de \LaTeX, pero esta plantilla dispone de comandos propios como por ejemplo el siguiente:

\imagen{escudoInfor}{Autómata para una expresión vacía}{.5}



\section{Listas de items}

Existen tres posibilidades:

\begin{itemize}
	\item primer item.
	\item segundo item.
\end{itemize}

\begin{enumerate}
	\item primer item.
	\item segundo item.
\end{enumerate}

\begin{description}
	\item[Primer item] más información sobre el primer item.
	\item[Segundo item] más información sobre el segundo item.
\end{description}
	
\begin{itemize}
\item 
\end{itemize}

\section{Tablas}

Igualmente se pueden usar los comandos específicos de \LaTeX o bien usar alguno de los comandos de la plantilla.

\tablaSmall{Herramientas y tecnologías utilizadas en cada parte del proyecto}{l c c c c}{herramientasportipodeuso}
{ \multicolumn{1}{l}{Herramientas} & App AngularJS & API REST & BD & Memoria \\}{ 
HTML5 & X & & &\\
CSS3 & X & & &\\
BOOTSTRAP & X & & &\\
JavaScript & X & & &\\
AngularJS & X & & &\\
Bower & X & & &\\
PHP & & X & &\\
Karma + Jasmine & X & & &\\
Slim framework & & X & &\\
Idiorm & & X & &\\
Composer & & X & &\\
JSON & X & X & &\\
PhpStorm & X & X & &\\
MySQL & & & X &\\
PhpMyAdmin & & & X &\\
Git + BitBucket & X & X & X & X\\
Mik\TeX{} & & & & X\\
\TeX{}Maker & & & & X\\
Astah & & & & X\\
Balsamiq Mockups & X & & &\\
VersionOne & X & X & X & X\\
} 

\capitulo{4}{Técnicas y herramientas}

\begin{comment}
Esta parte de la memoria tiene como objetivo presentar las técnicas metodológicas y las herramientas de desarrollo que se han utilizado para llevar a cabo el proyecto. Si se han estudiado diferentes alternativas de metodologías, herramientas, bibliotecas se puede hacer un resumen de los aspectos más destacados de cada alternativa, incluyendo comparativas entre las distintas opciones y una justificación de las elecciones realizadas. 
No se pretende que este apartado se convierta en un capítulo de un libro dedicado a cada una de las alternativas, sino comentar los aspectos más destacados de cada opción, con un repaso somero a los fundamentos esenciales y referencias bibliográficas para que el lector pueda ampliar su conocimiento sobre el tema.
\end{comment}

XAMPP es una distribución gratuita y de código abierto que facilita la instalación y el uso de un entorno de desarrollo web en un sistema local. Su nombre es un acrónimo que representa los componentes principales que incluye: 

\begin{itemize}
    \item \textbf{X} (Cross-platform): Indica que es multiplataforma, compatible con Windows, Linux, y macOS \cite{xampp}.
    \item \textbf{A} (Apache): Incluye el servidor web Apache, que es uno de los servidores web más populares del mundo.
    \item \textbf{M} (MariaDB/MySQL): Proporciona el sistema de gestión de bases de datos MariaDB (anteriormente MySQL).
    \item \textbf{P} (PHP): Incluye el lenguaje de scripting PHP, que es ampliamente utilizado para el desarrollo web del lado del servidor.
    \item \textbf{P} (Perl): También incluye Perl, otro lenguaje de programación.
\end{itemize}

XAMPP simplifica la instalación y configuración de estos componentes, lo que lo hace ideal para desarrolladores que desean crear, probar y depurar aplicaciones web localmente antes de desplegarlas en un servidor en producción \cite{xamppdocs}.

MySQL es un sistema de gestión de bases de datos relacional (RDBMS) que utiliza el lenguaje SQL (Structured Query Language) para gestionar y manipular datos. Es uno de los sistemas de bases de datos más utilizados en el mundo, especialmente en aplicaciones web \cite{mysql}. Algunas características clave de MySQL incluyen:

\begin{itemize}
    \item \textbf{Alto Rendimiento}: Optimizado para rendimiento y escalabilidad, capaz de manejar grandes cantidades de datos y transacciones.
    \item \textbf{Confiabilidad}: Proporciona características avanzadas como recuperación ante fallos y replicación para garantizar la integridad y disponibilidad de los datos.
    \item \textbf{Flexibilidad}: Soporta múltiples motores de almacenamiento (InnoDB, MyISAM, etc.) que permiten elegir el más adecuado según las necesidades específicas.
    \item \textbf{Compatibilidad}: Es compatible con múltiples plataformas y puede integrarse con una amplia variedad de lenguajes de programación y herramientas de desarrollo.
\end{itemize}

MySQL es utilizado por muchas aplicaciones y sitios web de alto tráfico, como Facebook, Twitter, y YouTube, debido a su eficiencia y robustez. Es una elección común para desarrolladores que buscan una solución de base de datos relacional confiable y de alto rendimiento \cite{mysqlcookbook}.

\capitulo{5}{Aspectos relevantes del desarrollo del proyecto}

\section{Etapas del proyecto}\label{etapas-del-proyecto}

\subsection{Investigación}
Durante la primera etapa del proyecto, se investigó acerca de \textit{FIWARE} y sus distintas aplicaciones, y se decidió enfocar la aplicación a usar el modelo información \textit{NGSI-LD} \ref{modelo-ngsi-ld}. 
Esta decisión fue impulsada por el reciente crecimiento de la plataforma \textit{FIWARE} y el modelo \textit{NGSI-LD}, gracias a la colaboración de las empresas más significativas en el desarrollo \textit{IOT}, como Telefónica \footnote{Proyecto de Telefónica de \textit{fiware-orion}, el cual lleva 6 años en marcha y a día de hoy se sigue ampliando. \url{https://github.com/telefonicaid/fiware-orion/tree/master}}.


\subsection{Despliegue del \textit{Orion Context Broker} y \textit{MongoDB}}

Tras un proceso de búsqueda e investigación de \textit{OpenData} reales que usaran el modelo \textit{NGSI-LD}, por ser este un estándar de uso libre, se descubrió que el número de conjuntos de datos que siguiesen este modelo de información era muy excaso aún. Por lo que se decidió crear un \textit{mock server} que usara NGSI-LD, en lugar de los distintos formatos que están en uso a día de hoy (\textit{JSON}, \textit{GEOJSON}, \textit{CSV}, \textit{XML}...), estandarizando así el uso de un mismo formato para todos los conjuntos de datos. Esto favorecería el desarrollo y mantenimiento de todas las aplicaciones que consuman este tipo de datos, no solo para las \textit{Smart Cities} sino en muchos otro ámbitos como \textit{Industria 4.0}.

Se comenzó con la instalación del \textit{Orion Context Broker} para la gestión de la información en tiempo real y, \textit{MongoDB}, que actúa como la base de datos subyacente para almacenar esta información de contexto. Todo ello se desplegó en un \textit{Docker}.
\imagen{OCB-Running}{\textit{Orion Context Broker} y la base de datos \textit{MongoDB} en funcionamiento}{.7}

Posteriormente se elaboró un \textit{script} que permitiese al usuario gestionarlo de manera más rápida, ya que, el \textit{Orion Context Broker} se gestiona mediante comandos, y para una mayor reusabilidad y mejor mantenimiento.
Con este mini programa se permite al usuario:
\begin{itemize}
    \item Iniciar y para el servidor, es decir encender el \textit{docker} de del \textit{Orion Context Broker} y el de la \textit{MongoDB}.
    \item Mostrar y borrar las entidades almacenadas.
    \item Iniciar y para el simulador. Explicaré su funcionamiento más adelante en \ref{simulador}.
    \item Instalar y desinstalar los \textit{dockers}.
\end{itemize}

\imagen{GestorOCB}{Programa gestor del OCB}{.8}

En esta etapa mediante la herramienta Postman\footnote{Herramienta para probar y gestionar APIs}, se hicieron pruebas con el Orion Context Broker, creando, consultando, actualizando y borrando entidades el servidor. Así se podía verificar que el servidor funcionaba correctamente y estaba preparado para almacenar las distintas entidades.
Las cuales posteriormente serían consultadas, actualizadas o eliminadas desde la aplicación con peticiones \textit{GET}\footnote{Se utiliza para solicitar datos de un recurso específico.}, \textit{PATCH}\footnote{Se utiliza para hacer actualizaciones parciales a un recurso existente.} o \textit{DELETE}\footnote{Se utiliza para eliminar un recurso específico.} de \textit{HTTP} y poder consultadas con peticiones \textit{GET} de \textit{HTTP} en el puerto 1026, que es el usado por el servidor.
\imagen{POST - Postman}{Ejemplo de uso de POST en Postman}{.9}\label{fig:post-postman}

\subsection{Simulador de conjunto de datos usando el modelo \textit{NGSI-LD}}\label{simulador}

Con las primeras cargas de datos se pudo observar que el modelo información \textit{NGSI-LD} todavía no está siendo muy usado por las diferentes \textit{APIs} de datos, por lo que se decidió crear un generador de datos que usara este modelo. Este es usado a modo de simulador, el cual genera de un conjunto de datos de \textit{parkings} de la ciudad de Burgos simulando recibir datos de sensores instalados en dichos \textit{parkings} y lo adapta al modelo \textit{NGSI-LD}. 

\imagen{GET - Postman}{Ejemplo de uso de GET en Postman}{.9}\label{fig:get-postman}


Con esto se pudo probar que el \textit{OCB}\footnotemark{} estaba funcionando correctamente y pudiendo así leer y escribir datos de manera correcta. También se introdujeron algunas entidades diferentes, no solo \textit{parkings}, para probar que los distintos conjuntos de datos podían convivir sin problema en el mismo\textit{OCB}\footnotemark[\value{footnote}]. 

\footnotetext{Orion Context Broker}

Para un mejor entendimiento de la estructura del proyecto, se adjunta una imagen con un esquema de la estructura \ref{fig:Estructura}.

\imagen{Estructura}{Estructura del proyecto}{.9}

\subsection{Desarrollo de la aplicación}

Como el objetivo del proyecto es realizar una aplicación funcional y, el número de conjuntos de datos disponibles con el modelo de información\textit{NGSI-LD} aún es mínimo, se comenzó con el desarrollo de la aplicación con los formatos tradicionales que se están usando a día de hoy (\textit{JSON}, \textit{GEOJSON}, \textit{CSV}, \textit{XML}...) y se dejó el \textit{OCB}\footnotemark[\value{footnote}] preparado para que cuando el número de conjuntos que utilicen el modelo de información\textit{NGSI-LD} sea mayor, se pueda utilizar.

Para el desarrollo de la aplicación se hizo un análisis de los posibles perfiles de usuarios que tendría y sus necesidades. Se clasificaron en dos grupos:

\begin{itemize}
    \item \textbf{Usuario estándar:} usaría la aplicación para acceder a un conjunto de datos en concreto en un momento puntual. Su principal objetivo sería obtener información acerca de un conjunto de manera rápida y precisa. Por ejemplo, un usuario quiere decidir a qué biblioteca de su ciudad acudir, por ello, consulta en conjunto y analiza las posibles opciones en función de los parámetros que más le interesen (aforo, horario, proximidad...).
    \item \textbf{Usuario avanzado:} este usuario sería un perfil más próximo a un analista de datos, que desearía consultar uno o varios conjuntos, durante un periodo de tiempo determinado con una frecuencia establecida.
\end{itemize}

Una vez establecidos los perfiles y las necesidades a cubrir, se empezó con el análisis y diseño de la aplicación. Y posteriormente su desarrollo, el cuál, ha sido complejo ya que requería de cubrir las necesidades de los usuarios actuales y de usuarios futuros con un modelo de información diferente, el modelo NGSI-LD, mencionado en el apartado \ref{modelo-ngsi-ld}.

Para su desarrollo se ha tenido en cuenta que entre los objetivos generales, se encuentren objetivos tales como permitir al usuario tener una experiencia sencilla, cómoda y rápida. Es por ello que un usuario puede marcar que parámetros desea visualizar de un conjunto, o marcar varios conjuntos como favoritos y poder acceder desde el menú de inicio de manera rápida.

Las funcionalidades principales de dicha aplicación, basándose en los perfiles de los usuarios potenciales y en los objetivos y necesidades a cubrir son:

\begin{itemize}
    \item \textbf{Consulta de conjuntos de datos en tiempo real:} esta sería la principal funcionalidad de la aplicación, permite al usuario consultar cualquier conjunto de datos de los que se encuentren incorporados en la aplicación. \imagen{puntos_de_carga_Valencia}{Ejemplo de consulta de datos en tiempo real}{.9} 
    \item \textbf{Añadir conjuntos de datos:} esta funcionalidad permite a los usuarios añadir nuevos conjuntos de datos, los cuales estarán visibles al resto de usuarios.
    \item \textbf{Añadir registros:} cualquier usuario puede comenzar a registrar un conjunto de datos con una frecuencia determinada, si la frecuencia es mayor con la que se actualizan en la \textit{API} de datos de donde se extraen, se notifica al usuario para que corrija su elección de periodicidad.
    \item \textbf{Gestionar registros:} un usuario que tenga uno más o registros guardándose para él, puede descargar o eliminar lo registros que desee.
    \item \textbf{Favoritos:} a la hora de consultar un conjunto de datos, un usuario puede marcar o desmarcar de favoritos un conjunto, esto modificará su ventana de \textit{home}, donde se muestran los conjuntos favoritos del usuario para poder acceder a ellos de una manera más rápida.
    \item \textbf{Personalizar vista de conjuntos de datos:} en los mapas, el usuario puede modificar los \textit{tooltips}, para poder mostrar los parámetros que él desee. También puede ordenar la tabla por orden alfabético según el campo que desee.
    \imagen{PCval2}{Ejemplo de consulta de datos en tiempo real 2}{.9}
\end{itemize}

Otras funcionalidades principales que solo tienen los administradores son:
\begin{itemize}
    \item \textbf{Gestionar conjuntos:} un administrador puede eliminar conjuntos. Esto permite que en caso de que un administrador detecte o se le comunique que alguno de los enlaces está roto o los datos están corruptos, podrá quitar ese conjunto de la base de datos.
    \item \textbf{Gestionar usuarios:} un administrador puede consultar todos los usuarios que se encuentran en la aplicación, cambiar el rol de algún usuario o eliminarlo.
    \item \textbf{Gestionar traducciones:} un puede añadir traducciones que se emplean para hacer que la lectura de datos por parte del usuario sea más cómoda y legible. Por ejemplo, un atributo que sea "PlLibr" en un conjunto de datos, se vea traducido como "Plazas libres"
\end{itemize}

\subsection{Revisión de código}
CUANDO SE TERMINE EL USO DE SONARQUBE

\begin{comment}

Este apartado pretende recoger los aspectos más interesantes del desarrollo del proyecto, comentados por los autores del mismo.
Debe incluir desde la exposición del ciclo de vida utilizado, hasta los detalles de mayor relevancia de las fases de análisis, diseño e implementación.
Se busca que no sea una mera operación de copiar y pegar diagramas y extractos del código fuente, sino que realmente se justifiquen los caminos de solución que se han tomado, especialmente aquellos que no sean triviales.
Puede ser el lugar más adecuado para documentar los aspectos más interesantes del diseño y de la implementación, con un mayor hincapié en aspectos tales como el tipo de arquitectura elegido, los índices de las tablas de la base de datos, normalización y desnormalización, distribución en ficheros3, reglas de negocio dentro de las bases de datos (EDVHV GH GDWRV DFWLYDV), aspectos de desarrollo relacionados con el WWW...
Este apartado, debe convertirse en el resumen de la experiencia práctica del proyecto, y por sí mismo justifica que la memoria se convierta en un documento útil, fuente de referencia para los autores, los tutores y futuros alumnos.
\end{comment}
\capitulo{6}{Trabajos relacionados}

\begin{comment}
Este apartado sería parecido a un estado del arte de una tesis o tesina. En un trabajo final grado no parece obligada su presencia, aunque se puede dejar a juicio del tutor el incluir un pequeño resumen comentado de los trabajos y proyectos ya realizados en el campo del proyecto en curso. 
\end{comment}

En este apartado quiero mencionar algunos trabajos relacionado con el proyecto, los cuales han proporcionado información de calidad al proyecto.

 \section{TFM "Desarrollo de un entorno de visualización de datos de smart-city basados en NGSI” de Mario Núñez Callejo}
 
 Este trabajo contribuye a la evolución de las ciudades inteligentes al proporcionar una herramienta que facilita la interpretación y el uso de datos urbanos complejos, apoyando así la toma de decisiones informadas y la gestión eficiente de los recursos urbanos.

En este proyecto también se elabora una herramienta para la interpretación de datos NGSI, solo que solo abarca datos de Cantabria, lo cual limita los usos de la herramienta, aunque se comparten algunos objetivos comunes como mejorar la interoperabilidad de los datos de ciudades inteligentes mediante la estandarización y el uso de NGSI.

\url{https://repositorio.unican.es/xmlui/bitstream/handle/10902/21941/434715.pdf}
\capitulo{7}{Conclusiones y Líneas de trabajo futuras}
\begin{comment}
Todo proyecto debe incluir las conclusiones que se derivan de su desarrollo. Éstas pueden ser de diferente índole, dependiendo de la tipología del proyecto, pero normalmente van a estar presentes un conjunto de conclusiones relacionadas con los resultados del proyecto y un conjunto de conclusiones técnicas. 
Además, resulta muy útil realizar un informe crítico indicando cómo se puede mejorar el proyecto, o cómo se puede continuar trabajando en la línea del proyecto realizado. 
\end{comment}

\section{Conclusiones}

Tras finalizar el proyecto se han obtenido diferentes conclusiones, las cuales he clasificado en tres grupos.
\begin{itemize}
    \item \textbf{Conclusiones del proyecto}
    \item \textbf{Conclusiones técnicas}
    \item \textbf{Conclusiones personales}
\end{itemize}

\subsection{Conclusiones del proyecto}
El proyecto ha concluido con éxito, cumpliendo todos y cada uno de los objetivos marcados al comienzo del mismo. Se ha conseguido realizar una aplicación funcional, cubriendo las necesidades de los usuarios y un costo de recursos mínimo.

\subsection{Conclusiones técnicas}
\begin{itemize}
    \item La aplicación es robusta, tiene un tratamiento de excepciones preparado para cubrir todos los errores. 
    \item Es escalable y fácilmente mantenible. Permite añadir nuevos formatos de manera rápida y simple.
    \item Tiene una interfaz amigable con el usuario, permitiendo que pueda personalizar los conjuntos que desea consultar y añadir nuevos. 
    \item Posee cifrado para que las contraseñas estén siempre seguras.
    \item Su implementación y uso es simple.
\end{itemize}

\subsection{Conclusiones personales}


\section{Posibles mejoras futuras}

% Bibliografía
\bibliographystyle{plain}  % Estilo de la bibiografía
\bibliography{bibliografia}

\end{document}
